% WS1 report fragment (issue #36): introduction, dataset and data preparation, experimental design.
% No preamble. The report (docs/report/main.tex, WS3) pulls this file in with \input and then places
% \wsoneIntro (Section 2), \wsoneDataset (Section 4), \wsoneDesign (Section 7) and
% \wsoneContribution (Section 11). Every number carries a "% source:" comment; "% TODO" marks a
% number that does not exist yet.
%
% Needs the packages booktabs and hyperref (the report preamble loads them).

\newcommand{\wsoneIntro}{%
\section{Introduction}

Stack Overflow answers millions of programming questions, and many are asked again in other words.
Moderators close such questions as duplicates of an older one, but only after someone notices the
match. Finding likely duplicates when a question is posted would save askers and answerers the same
work twice. The difficulty is scale: comparing every new question with every old one costs a number
of comparisons that grows with the square of the data.

We study one system for this task: Jaccard similarity of question text, estimated with MinHash
signatures, and found without comparing all pairs by locality-sensitive hashing (LSH), implemented
by us in Spark. Each stage of the search (how text becomes sets, how long the signatures are, how
LSH is set, how candidates are checked, what happens to very large buckets, how tags can narrow the
search) is a design choice that one experiment settles (E1 to E8). The project answers four questions.

\begin{description}
  \item[Q1 (accuracy).] How many known duplicates does our system find, and how close does it come
        to an exact search that checks every pair?
  \item[Q2 (theory).] Do MinHash and LSH behave as the course theory predicts on real questions, and
        which settings balance missed duplicates against wasted work?
  \item[Q3 (scale).] What slows the system down as it grows (oversized buckets, data moved between
        machines, memory), and how does run time change with more cores, a second machine and more data?
  \item[Q4 (secondary).] Does grouping questions by frequent tag sets, found with the SON algorithm,
        save work without losing many duplicates?
\end{description}

Our contributions are our own implementations of shingling, MinHash, LSH and SON; a streaming ETL for
the 69~GB data dump; an evaluation against the duplicates that moderators have labelled; and a
measured study of scale on two laptops run as a Spark cluster.
}

\newcommand{\wsoneDataset}{%
\section{Dataset and data preparation}

\paragraph{Source.} We use the Stack Exchange data dump, release 2026-09-30, from the Internet
Archive (licence CC BY-SA 4.0). Stack Overflow is the target: its archive is 68,929,051,847 bytes
% source: scripts/cluster/verify_archive.py (EXPECTED_ARCHIVE_SIZE)
and we keep the questions tagged \texttt{python}.
% TODO(#3): number of python questions and posts of Stack Overflow, from results/dataset_stats.md
Super User, a small site in the same format, is used for development on the two lighter laptops.

\paragraph{Super User.} The site has 516,669 questions and 753,503 answers, asked from 2008 to 2026
% source: results/dataset_stats.md (Questions, Answers, Year table)
(at most 47,106 questions in a year, in 2013).
% source: results/dataset_stats.md (Year table)
There are 90,843 links between posts, 14,641 of them marked as duplicates by moderators.
% source: results/dataset_stats.md (Duplicate links, Linked links)
Questions are 970 characters long on average; 35.4\% contain code, and 17,216 questions (3.3\%)
% source: results/dataset_stats.md (Mean body length, Bodies with code, Questions without an owner)
have no owner because the account was deleted. No question body is empty, and 3,212 questions
% source: results/dataset_stats.md (Empty bodies, `python` questions)
carry the \texttt{python} tag.

\paragraph{Preparation.} The raw data is XML inside one archive that must not be unpacked. Our ETL
streams \texttt{Posts.xml} and \texttt{PostLinks.xml} out of the archive, reads them row by row with
constant memory, and writes three Parquet tables partitioned by year: \texttt{posts\_meta} (all
questions and answers without text), \texttt{posts\_text} (question title and body, HTML kept) and
\texttt{post\_links}. The whole Super User archive takes 51 seconds and gives 1,270,172 posts.
% source: src/etl/README.md (Results on Super User); manifest of the run
HTML is removed and the text is cut into word shingles later, in the similarity stage, so that
experiment E1 can compare several choices.

\paragraph{Ground truth.} The answer key is the set of duplicate links whose two questions are both
in the data, with the newer question first. On Super User the 14,641 duplicate links give 14,639 pairs:
% source: src/etl/README.md (Results on Super User)
849 links (5.8\%) pointed from the older to the newer question and were turned around, and one link
between questions with the same timestamp was dropped; none referred to a missing question.

\paragraph{Limitations.} The answer key is incomplete: many true duplicates were never labelled, so
the precision we measure is a lower bound, and the recall we measure counts only duplicates that a
moderator found. Duplicate decisions are made by people and are not always consistent. Results on
Super User need not carry over to Stack Overflow, which has a different mix of topics and far more
code. Restricting Stack Overflow to \texttt{python} questions keeps the cluster run feasible but
leaves duplicates of python questions that are tagged otherwise out of the search.
}

\newcommand{\wsoneDesign}{%
\section{Experimental design}

\paragraph{Questions evaluated.} For every duplicate question in the answer key, the system returns a
ranked list of earlier questions (candidates). We count the question as found at rank $k$ if its
labelled original is among the first $k$ candidates. Questions of the answer key with no candidate
count as misses. Ties between equal similarities are broken by question id, so every run gives the same
ranking.

\paragraph{Metrics.} We report Recall@$K$ for $K = 1, 5, 10$, Precision@10, mean average precision
(MAP), NDCG@10 with one relevant item per question, and the reduction ratio
$1 - c / \binom{n}{2}$, the share of all pairs that LSH did not have to consider ($c$ candidate pairs among
$n$ questions). Precision is a lower bound for the reason given above. Every run adds one row to a
single table, \texttt{results/metrics.csv}, with its experiment, its settings and these values.

\paragraph{Baselines.} Two baselines show what LSH gives up and what it saves: an exact all-pairs
Jaccard search on nested samples of 10,000, 50,000 and, if time allows, 100,000 questions (E5), and
Spark's built-in \texttt{MinHashLSH} on the same samples. The exact search also provides the true
similarity of every pair, which experiment E3 uses to compare the measured chance of becoming a
candidate with the theoretical curve $1 - (1 - s^r)^b$.

\paragraph{Parameters.} Signature length, bands and rows, shingle size, the bucket strategy and the
random seed of every sample are settings in \texttt{configs/}, so each result can be reproduced from
its file. Scale is measured with the Spark event logs: total time, time per stage, shuffle volume,
peak memory and the slowest task against the median (E8).

\paragraph{Hardware and software.} Development and tests run in Spark local mode on Super User on two
lighter laptops. Full-scale runs use two laptops as a standalone Spark cluster with HDFS (replication 2)
under WSL2: an Intel Core i7-13650HX with 16~GB (master, NameNode, worker) and an AMD Ryzen 7 5800H with
% source: docs/team/jubayerkhan.md; docs/cluster_setup.md; README.md (team table)
16~GB (worker, DataNode). Software: Python 3.10 or newer, Spark 3.5.3 on Java 17, Hadoop 3.3, PyArrow
% source: requirements.txt (pyspark==3.5.3); README.md (Getting started)
for Parquet, and \texttt{7z} for the archive.
}

\newcommand{\wsoneContribution}{%
Satyam Bhanot wrote the data foundation and the evaluation: the streaming ETL that turns the data dump
into Parquet, the dataset profile, the answer key and the small sample files for the team, the evaluation
library and its command line, the error analysis, the scaling test scripts and the event-log parser, and
the repository conventions, contracts and automated checks (linter, tests in continuous integration, README
and submission checks). He also wrote this report's introduction, dataset and experimental-design sections.
}
